\documentclass[11pt,a4paper]{article}
\usepackage{amsmath,amssymb,amsfonts,amsthm}
\usepackage[margin=1in]{geometry}
\usepackage{booktabs}
\usepackage{graphicx}
\usepackage[font=small,labelfont=bf]{caption}
\usepackage{microtype}
\usepackage[colorlinks=true,linkcolor=black,citecolor=black,urlcolor=blue]{hyperref}

\newtheorem{prop}{Proposition}
\newtheorem{lem}[prop]{Lemma}
\newtheorem{rem}[prop]{Remark}
\newtheorem{defn}[prop]{Definition}

\def\R{\mathbb{R}}
\def\N{\mathbb{N}}
\def\E{\mathbb{E}}
\def\P{\mathbb{P}}
\def\xib{\xi}

\title{Bounded-support random bridges for mask-domain speech enhancement}
\author{}
\date{Draft scaffold --- \today}

\begin{document}
\maketitle

\begin{abstract}
\noindent
Generative speech enhancement generalises better than discriminative
enhancement to unseen acoustic conditions, but is held back by two documented
costs: it needs tens of network evaluations per utterance, and it can
hallucinate content that was not present in the input. We observe that both
costs follow from a modelling choice rather than from generative modelling
itself. Every diffusion and bridge model currently used for speech takes a
Gaussian driving process, whose state space is all of $\R^n$; the quantity being
estimated, a time--frequency mask, lives on $[0,1]$. We construct a random
bridge driven by a Wright--Fisher process, whose state space is $[0,1]$ by
construction, and derive its transition kernel as an exact Doob $h$-transform of
a finite binomial chain, avoiding the series expansion the Wright--Fisher
diffusion would otherwise require. On the empirical distribution of ideal ratio
masks the resulting bridge attains a given transport accuracy in fewer sampling
steps than a Gaussian bridge, by a factor of $2.3$ at five steps. Because a mask
in $[0,1]$ can only attenuate the observed signal, a mask-domain generative
model cannot synthesise content absent from its input, so the hallucination
failure mode is excluded by the construction rather than mitigated after the
fact. \emph{Sections~\ref{sec:exp} onward are a plan; results reported here are
one-dimensional and use exact conditional expectations.}
\end{abstract}

\section{Introduction}

Discriminative speech enhancement maps a noisy observation to a clean estimate
directly. Generative enhancement instead models the distribution of clean
speech and conditions on the observation. The generative route is more robust
when test conditions differ from training conditions
\cite{welker2022,lemercier2023}, which is the reason to pay for it, but two
costs are consistently reported. Diffusion-based systems require tens of network
evaluations per utterance, placing them one to two orders of magnitude above
discriminative systems in arithmetic cost \cite{comparison2026}. They can also
insert or substitute phonetic content, produce spurious breathing, and introduce
tonal artefacts; these hallucinations are described in the literature as an
open problem \cite{comparison2026,artifree2025}.

We attribute both costs to the choice of driving process rather than to
generative modelling itself. Existing diffusion and
bridge formulations for speech --- score-based models, Schr\"odinger bridges and
their one-step variants --- all drive the transport with Brownian motion, whose
state space is $\R^n$. When the estimated quantity is a time--frequency mask,
the state space is $[0,1]$, and a Gaussian transport spends part of every
trajectory outside it.

\paragraph{Contributions.}
\begin{enumerate}
\item We measure how much of a Gaussian bridge's trajectory leaves the feasible
set on real speech, for three representations, and show the violation is
substantial for masks and absent for the log-magnitude representation that
existing systems use (Section~\ref{sec:premise}).
\item We separate two properties that a support-matched driver may carry,
boundedness and monotonicity, and show by ablation that monotonicity is the
expensive one. This rules out the subordinator families and identifies the
requirement: bounded and non-monotone (Section~\ref{sec:which}).
\item We construct the Wright--Fisher bridge, deriving its transition kernel as
an exact $h$-transform of a finite binomial chain and so avoiding the infinite
mixture representation the diffusion limit requires
(Section~\ref{sec:wf}).
\item We show on the empirical mask distribution that the Wright--Fisher bridge
reaches a given accuracy in fewer sampling steps than a Gaussian bridge, and
explain the mechanism (Section~\ref{sec:why}).
\item We set out why a mask-domain generative model cannot hallucinate, and
specify the enhancement experiment that would test both claims
(Section~\ref{sec:exp}).
\end{enumerate}

\section{Background}

\subsection{Random bridges}

We use the formulation of \cite{companion}. For a c\`adl\`ag driving process
$\{Z_t\}_{t\in[0,T]}$ and laws $\Phi$, $\Psi$ with $\Phi\ll\mu_{Z_0}$ and
$\Psi\ll\mu_{Z_T}$, a process $\{\xib_t\}$ is a $(\Phi,\Psi)$-bridge under $Z$ if
its endpoints are distributed as a coupling of $\Phi$ and $\Psi$ and, conditional
on those endpoints, its finite-dimensional distributions agree with those of $Z$
conditioned likewise. The network estimates the $L^2$ projection
$f_\theta(\xib_s,s)\approx\E[Y\mid\xib_s]$ under squared error, and generation
integrates a transport whose terminal law is $\Psi$.

The definition constrains the driver only to be c\`adl\`ag. Its Gaussian
specialisation recovers diffusion bridges; the present paper uses a driver that
the definition admits but that no existing speech system employs.

\subsection{The Gaussian case}

With $Z_t=\sigma W_t$ and $t=sT$, writing $\gamma^2=\sigma^2T$, the bridge
marginal conditional on the endpoints is
\begin{equation}
\label{eq:gauss}
\xib_s=(1-s)x+sy+\gamma\sqrt{s(1-s)}\,\varepsilon,
\qquad \varepsilon\sim\mathcal{N}(0,I),
\end{equation}
and the exact transition from $s$ to $u$ pinned at an endpoint estimate
$\hat y$ is Gaussian with mean $\xib_s+\frac{u-s}{1-s}(\hat y-\xib_s)$ and
variance $\gamma^2\frac{(u-s)(1-u)}{1-s}$. Score-based enhancement, diffusion
bridges and Schr\"odinger bridges for speech differ in their schedule,
parameterisation and training objective, but all place the state in $\R^n$ and
all inherit \eqref{eq:gauss}'s unbounded support.

\section{Support violation in speech representations}
\label{sec:premise}

\paragraph{Measurement.} We take real speech, add noise at $5$\,dB, and form
three representations of the enhancement problem: the linear magnitude
spectrogram, with support $[0,\infty)$; the ideal ratio mask, with support
$[0,1]$; and the phase, which is circular. For each we construct the Gaussian
bridge \eqref{eq:gauss} between the true endpoints and record the fraction of
path mass outside the feasible set.

\begin{table}[h]
\centering\small
\caption{Fraction of a Gaussian bridge's trajectory outside the feasible set, on
$149$\,s of real speech. Excess is relative to each representation's own scale.
The log-magnitude row is a control: its support is unconstrained, and the
measurement returns zero.}
\label{tab:premise}
\begin{tabular}{llrr}
\toprule
Representation & Support & Outside & Mean excess \\
\midrule
Linear magnitude & $[0,\infty)$ & $11.1\%$ & $0.074$ \\
Ratio mask       & $[0,1]$      & $8.0\%$  & $0.063$ \\
Phase            & circular     & $1.3\%$  & $0.069$ \\
Log magnitude (control) & $\R$  & $0.00\%$ & $0$ \\
\bottomrule
\end{tabular}
\end{table}

\paragraph{The mask case.} The magnitude constraint is removed by taking logarithms, which is the
representation existing systems use and which the control row of
Table~\ref{tab:premise} confirms is unconstrained. No such transformation exists
for a mask. It is bounded on both sides, and the standard
device --- a sigmoid on the network output --- is a projection applied after the
fact, not a process that respects the constraint.

The mask distribution also concentrates at the boundary: $81\%$ of
time--frequency bins fall in the lowest decile and $77\%$ of mass lies within
$0.05$ of $0$ or $1$. A transport that overshoots does so into the region where
the target places most of its mass.

\section{Boundedness and monotonicity}
\label{sec:which}

A driver matched to $[0,1]$ may be bounded, monotone, or both. The gamma and
Poisson bridges derived in \cite{companion} are both: they are driven by
subordinators, so their paths cannot decrease. We separate the two properties
by ablation, taking a Gaussian bridge and forbidding its steps to decrease while
leaving the driver, the noise and the support handling unchanged.

\begin{table}[h]
\centering\small
\caption{Monotonicity is the expensive property. Targets are Gamma laws of
varying shape at fixed mean, so only the behaviour of the density at the
boundary changes; $1$-Wasserstein distance at $50$ steps, exact conditional
means.}
\label{tab:mono}
\begin{tabular}{lrrrr}
\toprule
Shape $k$ & Gaussian & Gaussian, monotone & Gamma bridge & monotone / free \\
\midrule
$0.5$ & $0.106$ & $1.412$ & $0.362$ & $13.3\times$ \\
$2.0$ & $0.033$ & $0.498$ & $0.083$ & $15.1\times$ \\
$8.0$ & $0.012$ & $0.102$ & $0.019$ & $8.3\times$ \\
\bottomrule
\end{tabular}
\end{table}

Forbidding decrease costs between $8$ and $15$ times, against the factor of two
to three by which the gamma bridge trails a free Gaussian. Subordinator-driven
bridges are therefore unsuited to a quantity that is not ordered relative to its
starting point. A mask is not so ordered: a given bin's ideal value may lie
above or below any initial estimate. The requirement is a driver that is bounded
and non-monotone.

\section{The Wright--Fisher bridge}
\label{sec:wf}

\subsection{The driver}

The Wright--Fisher diffusion on $[0,1]$,
\begin{equation}
\mathrm{d}X_t=\tfrac12\big(\theta_1(1-X_t)-\theta_2X_t\big)\mathrm{d}t
+\sqrt{X_t(1-X_t)}\,\mathrm{d}W_t,
\end{equation}
is the canonical diffusion with this state space. It is continuous, hence
c\`adl\`ag, so it is admitted by the bridge definition; its diffusion
coefficient is state dependent, so it has neither homogeneous nor independent
increments and the L\'evy results of \cite{companion} do not reach it. Its
paths are not monotone.

With $\theta_1=\theta_2=0$ both boundaries absorb, and the law of $X_T$ carries
atoms at $0$ and $1$ together with a density between. Admissibility depends on this. The requirement $\Psi\ll\mu_{Z_T}$ fails for a
mask distribution with boundary atoms when the boundaries are inaccessible,
which is the case under strong mutation. The neutral process is the member of
the family whose terminal law has the same atomic structure as the data.

\subsection{Transition kernel}

The transition density of the Wright--Fisher diffusion is available only as an
infinite mixture over a death process \cite{griffiths1979,jenkins2017}. The
finite population model of which it is the diffusion limit has no such
difficulty. With population size $M$, one generation is a binomial resampling,
\begin{equation}
\label{eq:wfchain}
P(i\to j)=\binom{M}{j}\Big(\tfrac{i}{M}\Big)^{j}\Big(1-\tfrac{i}{M}\Big)^{M-j},
\qquad i,j\in\{0,\dots,M\},
\end{equation}
an exact finite stochastic matrix, with $0$ and $M$ absorbing.

\begin{prop}[Wright--Fisher bridge kernel]
\label{prop:wf}
Let $\{\xib\}$ be the bridge under the chain \eqref{eq:wfchain} pinned at
$\xib_1=y$, and let $a$ and $b$ be the generation counts on $[s,u]$ and $[u,1]$.
Then
\begin{equation}
\P(\xib_u=j\mid\xib_s=i,\ \xib_1=y)
=\frac{P^{a}[i,j]\;P^{b}[j,y]}{P^{a+b}[i,y]},
\end{equation}
which is the Doob $h$-transform of \eqref{eq:wfchain} with space--time harmonic
function $h_u(j;y)=P^{b}[j,y]$.
\end{prop}

\begin{proof}
Immediate from the Markov property of the chain and Bayes' rule; the denominator
is the normalisation $\sum_j P^{a}[i,j]P^{b}[j,y]=P^{a+b}[i,y]$.
\end{proof}

Every quantity is a power of a finite matrix. For $M=128$ and a few hundred
generations the powers are precomputed once.

\begin{rem}[Absorbed states]
Because both boundaries absorb, a path that has reached $0$ or $1$ cannot reach
a pin that disagrees with it, and the kernel of Proposition~\ref{prop:wf} is
undefined there. Such paths remain where they are, which is what absorption
means. In sampling this arises when the endpoint estimate moves after a
trajectory has absorbed.
\end{rem}

\subsection{Sampling}

Training is unchanged from the Gaussian case: sample a bridge time, sample the
state from the marginal, regress the endpoint under squared error. Generation
alternates evaluating $f_\theta$ and sampling Proposition~\ref{prop:wf} with the
estimate in place of the unknown endpoint, exactly as ancestral diffusion
sampling substitutes a point estimate for the posterior.

\section{Sampling efficiency}
\label{sec:why}

\begin{table}[h]
\centering\small
\caption{Transport accuracy against sampling steps on the empirical distribution
of ideal ratio masks. $1$-Wasserstein distance, exact conditional means,
$M=128$, $256$ generations, $5\%$ relative error injected into the endpoint
estimate. Floor $0.0027$.}
\label{tab:wf}
\begin{tabular}{lrrrrr}
\toprule
Steps & $2$ & $5$ & $10$ & $20$ & $50$ \\
\midrule
Gaussian, projected at the end & $0.0715$ & $0.0485$ & $0.0335$ & $0.0253$ & $\mathbf{0.0179}$ \\
Gaussian, projected each step  & $0.0715$ & $0.0487$ & $0.0341$ & $0.0262$ & $0.0198$ \\
Wright--Fisher                 & $\mathbf{0.0561}$ & $\mathbf{0.0208}$ & $\mathbf{0.0192}$ & $\mathbf{0.0200}$ & $0.0207$ \\
\midrule
Ratio & $0.79$ & $0.43$ & $0.57$ & $0.79$ & $1.15$ \\
\bottomrule
\end{tabular}
\end{table}

The Wright--Fisher bridge is ahead below fifty steps and by a factor of $2.3$ at
five. The two cross at fifty.

\paragraph{Mechanism.} A Gaussian transport towards a boundary-concentrated
target overshoots, and correcting an overshoot costs steps: the drift
$\frac{\hat y-\xib_s}{1-s}$ must first reverse the excursion and then re-approach.
Given enough steps the correction completes, which is why the Gaussian catches
up at fifty. The Wright--Fisher transport cannot overshoot, because its
diffusion coefficient $\sqrt{x(1-x)}$ vanishes at both boundaries, so the steps
it does take are not spent undoing earlier ones. The advantage is largest at small step budgets, where generative enhancement is
currently impractical.

The crossover at fifty steps follows from the same account. An unconstrained
transport is the more flexible of the two given a sufficient budget, and the
constraint ceases to pay once the budget is ample.

\section{Mask-domain enhancement}
\label{sec:exp}

\subsection{Support of the estimate}

Let $\hat S=m\odot X$ with $m\in[0,1]$ and $X$ the observed noisy spectrogram.
Then $|\hat S_{tf}|\le|X_{tf}|$ for every bin, so the estimate is supported on
the observation. A mask-domain generative model can suppress or fail to suppress; it cannot place
energy where the input had none. The phonetic
insertions and spurious vocalisations reported for waveform and complex-spectrum
diffusion \cite{comparison2026,artifree2025} are not available to it.

The guarantee belongs to the output parameterisation rather than to the driver.
The driver supplies a distribution over masks that does not leave $[0,1]$, so
the guarantee is not obtained by projecting a model that violates it.

\subsection{Protocol}

\emph{To be run.} Enhancement on VoiceBank--DEMAND with a mismatched evaluation
set, since mismatch is where the generative route earns its cost. The transport
runs from a cheap deterministic mask computed from the observation to the ideal
ratio mask, both conditioned on the same noisy spectrogram; this uses the
bridge's arbitrary source distribution rather than starting from noise.

Arms: a discriminative mask predictor; a Gaussian mask bridge, which is the
ablation isolating the driver; a Wright--Fisher mask bridge; and a published
waveform or complex-spectrum diffusion system, which is where hallucination has
been reported. All share architecture, budget and evaluation code.

Measures: PESQ, ESTOI and SI-SDR for quality; network evaluations and wall-clock
for cost; and, for hallucination, word and phoneme error rate under a fixed ASR
system, on which mask-domain models are constrained by construction and waveform models
are not.

\subsection{Falsification}

If a Gaussian mask bridge matches the Wright--Fisher bridge at five to ten steps
on real data, the mechanism of Section~\ref{sec:why} does not survive contact
with a learned conditional expectation, and the contribution reduces to the
derivation. If mask-domain enhancement cannot reach the quality of
complex-spectrum systems at any step count, the hallucination guarantee is
purchased at a price the application will not pay.

\section{Limitations}

The results in Tables~\ref{tab:mono} and~\ref{tab:wf} are one-dimensional, use
exact conditional expectations rather than a trained network, and treat
time--frequency bins independently. They isolate a mechanism and measure its size in a setting where nothing else
can account for the difference. They do not establish that it survives at
spectrogram scale with a learned estimator. The speech
measurements of Table~\ref{tab:premise} use a small corpus and synthetic
additive noise. No enhancement experiment has been run.

\begin{thebibliography}{99}
\bibitem{companion} [Companion paper: random bridges as stochastic transports for generative models.]
\bibitem{welker2022} Welker, S., Richter, J., Gerkmann, T.
\emph{Speech enhancement with score-based generative models in the complex STFT domain}. Interspeech 2022.
\bibitem{lemercier2023} Lemercier, J.-M., Richter, J., Welker, S., Gerkmann, T.
\emph{StoRM: a diffusion-based stochastic regeneration model for speech enhancement and dereverberation}. IEEE/ACM TASLP 2023.
\bibitem{comparison2026} \emph{A comparison of generative and discriminative methods for speech enhancement: robustness, complexity, and hallucination}. arXiv:2606.02913.
\bibitem{artifree2025} \emph{ArtiFree: detecting and reducing generative artifacts in diffusion-based speech enhancement}. arXiv:2509.19495.
\bibitem{griffiths1979} Griffiths, R.C. \emph{A transition density expansion for a multi-allele diffusion model}. Adv. Appl. Prob. 1979.
\bibitem{jenkins2017} Jenkins, P.A., Span\`o, D. \emph{Exact simulation of the Wright--Fisher diffusion}. Ann. Appl. Prob. 2017.
\end{thebibliography}

\end{document}
